\documentclass[10pt,a4paper]{article}

\usepackage[utf8]{inputenc}
\usepackage[T1]{fontenc}
\usepackage[brazil]{babel}
\usepackage[margin=1.45cm]{geometry}
\usepackage{booktabs}
\usepackage{array}
\usepackage{enumitem}
\usepackage{listings}
\usepackage{xcolor}
\usepackage{hyperref}
\hypersetup{colorlinks=true, linkcolor=black, urlcolor=blue}

\setlength{\parindent}{0pt}
\setlength{\parskip}{3pt}
\setlist[itemize]{leftmargin=*, topsep=2pt, itemsep=1pt}
\setlist[enumerate]{leftmargin=*, topsep=2pt, itemsep=1pt}

\lstset{
  basicstyle=\ttfamily\small,
  breaklines=true,
  columns=fullflexible,
  showstringspaces=false,
  frame=single,
  framerule=0.2pt,
  rulecolor=\color{black!25},
  xleftmargin=2pt,
  xrightmargin=2pt
}

\title{\vspace{-1.1cm}\textbf{Apresentação e Defesa TR1}\\
\large roteiro completo de 10+ minutos, teoria e excertos de código}
\author{Gustavo Henrique Andrade Cavalcanti \and Fernando Augusto Hortencio}
\date{Teleinformática e Redes 1 - Universidade de Brasília - 2026}

\begin{document}
\maketitle
\vspace{-0.55cm}

\begin{abstract}
Este é o material único para apresentação e defesa do Simulador TR1. O
roteiro mira 11 a 12 minutos de fala e inclui divisão por integrante,
conceitos teóricos, excertos de código e perguntas prováveis. O relatório
formal permanece em \texttt{relatorio\_tr1.pdf}.
\end{abstract}

\section{Planejamento de tempo}

\begin{tabular}{p{4.1cm}p{1.8cm}p{3.1cm}p{5.5cm}}
\toprule
\textbf{Bloco} & \textbf{Tempo} & \textbf{Fala} & \textbf{Foco} \\
\midrule
Abertura e visão geral & 1 min & Fernando & problema, objetivo e fluxo completo \\
Conceitos antes do código & 2 min & Fernando/Gustavo & camadas, enlace, física e ruído \\
Arquitetura & 1 min 30 s & Fernando & threads, filas e meio \\
Camada de enlace & 3 min & Gustavo & quadros, edc, crc e hamming \\
Camada física e meio & 2 min 30 s & Gustavo & sinais, portadora, constelação e ruído \\
Interface e validação & 1 min & Fernando & gtk, testes e resultados \\
Fechamento & 30 s & Fernando & conclusão \\
\bottomrule
\end{tabular}

\section{Ajuste de discurso para a primeira fase}

Como todos os grupos receberam a mesma proposta, a apresentação não precisa
tratar o enunciado como novidade. A abertura deve ser curta e deve deslocar o
foco para as decisões da nossa entrega:

\begin{itemize}
  \item onde exatamente o ruído entra;
  \item como o bit vira amostra em volts;
  \item como o receptor decide o bit ou símbolo recebido;
  \item por que EDC vem antes do Hamming no transmissor;
  \item como os testes provam que o pipeline completo funciona.
\end{itemize}

Frase-guia: a proposta era comum a todos os grupos; o diferencial da nossa
entrega está no pipeline completo, manual e testável: texto, bits, quadros,
sinal, ruído, demodulação, verificação e texto recuperado.

\section{Abertura - Fernando}

Fala sugerida: como todos receberam a mesma proposta, a nossa apresentação vai
focar menos em repetir o enunciado e mais em mostrar como implementamos cada
decisão. A mensagem começa como texto, vira bits, é organizada em quadros,
recebe detecção e correção de erros, vira sinal em volts, atravessa um meio
com ruído gaussiano e depois é reconstruída no receptor.

\begin{center}
\texttt{texto -> bits -> enlace -> física -> meio ruidoso -> física -> enlace -> texto}
\end{center}

Pontos obrigatórios: a interface principal é gtk; tx e rx são separados; o
ruído entra no meio; crc, hamming, enquadramento e modulações são manuais; a
defesa deve insistir em decisões de projeto, não só listar protocolos.

\section{Conceitos antes do código}

\subsection{Visão de camadas - Fernando}

\begin{tabular}{p{2.7cm}p{5.2cm}p{5.1cm}}
\toprule
\textbf{Camada} & \textbf{Papel} & \textbf{Arquivo} \\
\midrule
Aplicação & texto <-> bits & \texttt{camada\_aplicacao.py} \\
Enlace & quadros, edc, hamming & \texttt{camada\_enlace.py} \\
Física & banda-base e portadora & \texttt{camada\_fisica.py} \\
Meio & ruído e potência & \texttt{meio\_comunicacao.py} \\
Simulação & tx, canal e rx & \texttt{simulador.py} \\
\bottomrule
\end{tabular}

\subsection{Enlace - Gustavo}

A camada de enlace resolve três problemas: definir fronteiras de quadro,
detectar alterações e, quando hamming está ligado, corrigir erro simples.
Os protocolos são: contagem, byte stuffing, bit stuffing, paridade, checksum,
crc-32 e hamming(8,4).

\subsection{Física - Gustavo}

A camada física transforma bits em amostras de tensão. Em banda-base, cada bit
vira níveis de tensão. Na portadora, grupos de bits escolhem amplitude,
frequência, fase ou ponto de constelação. QPSK usa 2 bits por símbolo; 16-QAM
usa 4 bits por símbolo, mas é mais sensível ao ruído.

\section{Arquitetura - Fernando}

\begin{lstlisting}[language=Python]
th_tx = threading.Thread(target=_rotina_tx, args=(config, fila_tx, resultados))
th_rx = threading.Thread(target=_rotina_rx, args=(config, fila_rx, resultados))

th_tx.start()
th_rx.start()

sinal_limpo = fila_tx.get()
sinal_ruidoso = meio_comunicacao.transmitir(
    sinal_limpo, config["ruido_media"], config["ruido_sigma"])
fila_rx.put(sinal_ruidoso)
\end{lstlisting}

Explicação: \texttt{th\_tx} desce a pilha no transmissor; \texttt{th\_rx}
sobe a pilha no receptor; \texttt{fila\_tx} leva o sinal limpo ao canal; o
meio soma ruído; \texttt{fila\_rx} entrega o sinal ruidoso ao receptor. Isso
mostra que tx e rx não se chamam diretamente.

\begin{lstlisting}[language=Python]
for byte in texto.encode("utf-8"):
    for i in range(7, -1, -1):
        bits.append((byte >> i) & 1)
\end{lstlisting}

Explicação: o texto vira bytes e cada byte é lido do bit mais significativo
ao menos significativo. Essa lista de bits é a entrada da camada de enlace.

\section{Camada de enlace - Gustavo}

\subsection{Pipeline}

\begin{lstlisting}[language=Python]
for i in range(0, len(bits), tam_bits):
    bloco = bits[i:i + tam_bits]
    bloco = ADICIONAR_EDC[config["deteccao"]](bloco)
    if config["correcao"] == "hamming":
        bloco = codificar_hamming(bloco)
    payloads.append(bloco)
return ENQUADRAR[config["enquadramento"]](payloads)
\end{lstlisting}

Explicação: o payload é dividido em blocos; o edc é anexado; hamming protege
payload + edc; por fim, o quadro é enquadrado. No receptor a ordem é inversa:
desenquadrar, corrigir e verificar.

\subsection{Enquadramento}

\begin{lstlisting}[language=Python]
if byte in (FLAG_BYTE, ESC_BYTE):
    quadro.append(ESC_BYTE)
quadro.append(byte)
\end{lstlisting}

No byte stuffing, \texttt{FLAG\_BYTE} delimita quadro e \texttt{ESC\_BYTE}
protege dados especiais.

\begin{lstlisting}[language=Python]
if uns == 5:
    trem.append(0)
    uns = 0
\end{lstlisting}

No bit stuffing, a flag é \texttt{01111110}. Depois de cinco bits 1 seguidos,
o transmissor insere 0 para impedir que dados imitem a flag.

\subsection{CRC-32}

\begin{lstlisting}[language=Python]
for bit in reversed(byte):
    if (crc ^ bit) & 1:
        crc = (crc >> 1) ^ POLI_CRC32_REFLETIDO
    else:
        crc >>= 1
return crc ^ 0xFFFFFFFF
\end{lstlisting}

Explicação: o crc é calculado bit a bit. O xor representa soma em gf(2), o
polinômio refletido é \texttt{0xedb88320}, e o vetor \texttt{123456789} deve
gerar \texttt{0xcbf43926}. Não há uso de \texttt{zlib}.

\subsection{Hamming(8,4)}

\begin{lstlisting}[language=Python]
s1 = (bloco[0] + bloco[2] + bloco[4] + bloco[6]) % 2
s2 = (bloco[1] + bloco[2] + bloco[5] + bloco[6]) % 2
s4 = (bloco[3] + bloco[4] + bloco[5] + bloco[6]) % 2
sindrome = s4 * 4 + s2 * 2 + s1
par_geral = sum(bloco) % 2
\end{lstlisting}

\begin{tabular}{p{3cm}p{3cm}p{7cm}}
\toprule
\textbf{Síndrome} & \textbf{Paridade geral} & \textbf{Resultado} \\
\midrule
0 & par & sem erro \\
diferente de 0 & ímpar & erro simples corrigível \\
0 & ímpar & erro no bit de paridade geral \\
diferente de 0 & par & erro duplo detectado \\
\bottomrule
\end{tabular}

Explicação: a síndrome aponta a posição do erro simples. A paridade geral
impede que dois erros sejam tratados como um erro simples. O custo é dobrar o
tamanho, porque 4 bits viram 8 bits.

\section{Camada física e meio - Gustavo}

\subsection{Banda-base}

Em banda-base, bits viram níveis de tensão. NRZ-Polar usa +V e -V; Manchester
usa transição no meio do bit; Bipolar AMI usa 0 V para zero e alterna a
polaridade dos bits 1.

Para o Manchester, a convenção adotada é a de Tanenbaum/G.E.~Thomas: o bit 0
faz transição baixo-para-alto e o bit 1 faz transição alto-para-baixo. A
demodulação compara a média da primeira metade do bit com a média da segunda
metade.

\subsection{Portadora e i/q}

Ponto importante para defesa: a modulação gera o sinal limpo. O ruído não faz
parte da modulação; ele é somado depois, no meio de comunicação, amostra por
amostra.

\begin{lstlisting}[language=Python]
def onda(i, q, ciclos):
    N = AMOSTRAS_POR_SIMBOLO
    amostras = []
    for n in range(N):
        angulo = 2 * math.pi * ciclos * n / N
        amostras.append(i * math.cos(angulo) - q * math.sin(angulo))
    return amostras
\end{lstlisting}

Explicação: cada símbolo é uma sequência de amostras. \texttt{i} multiplica o
cosseno e \texttt{q} multiplica o seno. Escolher \texttt{(i, q)} é escolher
amplitude e fase.

\begin{lstlisting}[language=Python]
i = 2 / N * soma_cos
q = -2 / N * soma_sen
return i, q
\end{lstlisting}

Explicação: o receptor estima as componentes i/q por correlação e depois
escolhe o ponto válido mais próximo da constelação.

\subsection{QPSK e 16-QAM}

QPSK transmite 2 bits por símbolo e usa mapeamento gray. 16-QAM transmite 4
bits por símbolo usando níveis \(-3\), \(-1\), \(1\), \(3\) em cada eixo. É
mais eficiente, mas os pontos ficam mais próximos, então o ruído afeta mais.

\subsection{Meio e potência}

\begin{lstlisting}[language=Python]
return [amostra + random.gauss(media, sigma) for amostra in sinal]
\end{lstlisting}

Explicação: cada amostra recebe ruído gaussiano independente. \texttt{media}
desloca o sinal; \texttt{sigma} controla a intensidade do ruído.

\begin{lstlisting}[language=Python]
return sum(amostra * amostra for amostra in sinal) / len(sinal)
\end{lstlisting}

Explicação: potência média é a média de \(v^2\), assumindo resistência de 1
ohm.

\section{Interface e validação - Fernando}

\begin{lstlisting}[language=Python]
resultado = simulador.executar_simulacao(config)
diagnostico = diagnosticar_camadas(
    resultado["tx_bits_aplicacao"], resultado, config)
\end{lstlisting}

Explicação: a interface não reimplementa protocolos. Ela chama a função real
do simulador, calcula diagnóstico por fase e mostra bits, sinais, métricas e
relatório de quadros.

\begin{lstlisting}[language=Python]
bits_123 = app.texto_para_bits("123456789")
ok &= testar("crc32 123456789 == 0xCBF43926",
             enlace.calcular_crc32(bits_123) == 0xCBF43926)
\end{lstlisting}

Explicação: esse teste valida o crc contra um vetor conhecido. A suíte também
cobre aplicação, enquadramentos, edc, hamming, modulações e simulação
completa.

\section{Fechamento - Fernando}

O projeto cumpre o enunciado porque separa transmissor e receptor, usa gtk
como interface principal, implementa enlace e física manualmente, usa sinal em
volts, injeta ruído gaussiano no meio e valida a transmissão com testes. A
principal contribuição é tornar visível o caminho completo da informação:
texto, bits, quadros, sinal, ruído e texto recuperado.

\section{Perguntas prováveis}

\textbf{1. Como vocês modulam com ruído?} Nós não modulamos ``com'' ruído.
Primeiro a camada física gera o sinal limpo a partir dos bits. Depois
\texttt{meio\_comunicacao.transmitir} soma \texttt{random.gauss(media, sigma)}
em cada amostra. O receptor demodula esse sinal ruidoso.

\textbf{2. Como foi feita a modulação?} Em banda-base, cada bit vira um
conjunto de amostras em volts: NRZ usa \(+V\) e \(-V\), Manchester usa
transição no meio do bit e Bipolar AMI alterna a polaridade dos bits 1. Na
portadora, cada símbolo vira uma onda \(I\cos - Q\sin\); ASK muda
amplitude, FSK muda frequência, QPSK usa 2 bits por símbolo e 16-QAM usa 4 bits
por símbolo.

\textbf{3. O ruído entra antes ou depois da modulação?} Depois. A ordem é:
bits \(\to\) modulação \(\to\) sinal limpo \(\to\) meio ruidoso \(\to\) sinal
recebido \(\to\) demodulação.

\textbf{4. Qual é o jeito certo de colocar o bit de paridade?} No código, a
paridade é calculada sobre o bloco inteiro e anexada como um byte de EDC: sete
zeros mais o bit de paridade. Ela não é colocada no final de cada byte original
da mensagem.

\textbf{5. Por que não usar zlib para crc?} Porque o enunciado proíbe
bibliotecas prontas para os protocolos. O crc foi implementado bit a bit.

\textbf{6. Por que usar paridade se CRC-32 é melhor?} Porque o trabalho pede
comparar mecanismos. A paridade é simples e tem baixo custo, mas só detecta
quantidade ímpar de erros. O CRC-32 detecta muito melhor, principalmente erros
em rajada, mas insere 32 bits por bloco.

\textbf{7. Por que hamming vem depois do edc?} Para proteger também os bits de
verificação anexados ao payload. No receptor, o Hamming tenta corrigir antes de
verificar o EDC.

\textbf{8. Por que hamming(8,4), e não hamming(7,4)?} O bit extra de paridade
geral permite detectar erro duplo e mantém alinhamento em bytes.

\textbf{9. Expliquem o Hamming do projeto.} Cada 4 bits de dados viram 8 bits:
três bits de paridade de Hamming, quatro bits de dados e um bit de paridade
geral. No receptor, a síndrome aponta erro simples corrigível; a paridade geral
permite diferenciar erro simples de erro duplo.

\textbf{10. Quais são as vantagens e desvantagens de NRZ-Polar e 16-QAM?}
NRZ-Polar é simples e robusto para simulação em banda-base, mas pode ter longas
sequências sem transição e componente DC. 16-QAM carrega 4 bits por símbolo e é
mais eficiente espectralmente, mas os pontos da constelação ficam próximos e a
demodulação fica mais sensível ao ruído.

\textbf{11. Por que 16-QAM é mais sensível a ruído?} Porque os pontos da
constelação ficam mais próximos. Uma variação menor já pode empurrar o símbolo
para a decisão errada.

\textbf{12. O que significa ``não dá para ver o sinal que está chegando'' na
portadora?} Dá para ver o sinal recebido como amostras ruidosas. O que não dá é
olhar a onda e ler diretamente os bits como em um desenho ideal. Por isso o
receptor estima \(I\) e \(Q\) por correlação e escolhe o ponto válido mais
próximo da constelação.

\textbf{13. O que acontece se o ruído aumentar muito?} O demodulador passa a
errar bits. O Hamming pode corrigir erro simples por bloco, mas erros múltiplos
podem ser apenas detectados ou até escapar dependendo do EDC.

\textbf{14. Como funciona o enquadramento com flags?} No byte stuffing, o
quadro começa e termina com \texttt{0x7E}; se \texttt{0x7E} ou \texttt{0x7D}
aparecem nos dados, o transmissor insere \texttt{0x7D} antes. No bit stuffing,
a flag é \texttt{01111110}; depois de cinco bits 1 seguidos no payload, o
transmissor insere um 0.

\textbf{15. Como vocês testaram?} \texttt{testes.py} cobre conversão
texto/bits, os três enquadramentos, paridade, checksum, CRC, Hamming,
modulações banda-base, modulações por portadora e simulação completa. O CRC
também é conferido com o vetor clássico \texttt{123456789 -> 0xCBF43926}.

\textbf{16. A interface web substitui a gtk?} Não. A interface principal é gtk; a
web é opcional e chama a mesma função.

\textbf{17. Como provar que o crc está correto?} Pelo vetor clássico
\texttt{123456789}, que deve retornar \texttt{0xcbf43926}.

\section{Checklist de domínio}

Fernando: \texttt{executar\_simulacao}, \texttt{\_rotina\_tx},
\texttt{\_rotina\_rx}, filas, interface gtk, diagnóstico e testes.

Gustavo: enquadramento, crc-32, hamming, \texttt{onda},
\texttt{correlacionar}, qpsk, 16-qam, ruído e potência média.

\end{document}
